\begin{exercisechunk}
\item \textbf{Construct a recurrent ReLU simulation of a finite automaton.}
\label[exercise]{ex:l14-recurrent-dfa}
\exerciseprofile{2}{2}{2}{\exproof\quad\exconstruction\quad\exinterpretation}{%
Construct one fixed neural transition cell and distinguish its parameter count
from the number of times it is executed.}
The Turing-machine model separates finite control from external memory.
This exercise constructs the finite-control part directly and shows how
recurrence reuses one transition network for every input position.

Let \(Q\) and \(\Sigma\) be finite sets, and let
\(\delta:Q\times\Sigma\to Q\) be the transition function of a deterministic
finite automaton. Encode \(q\in Q\) and \(a\in\Sigma\) by the one-hot vectors
\(e_q\in\{0,1\}^{|Q|}\) and
\(e_a\in\{0,1\}^{|\Sigma|}\).
\begin{enumerate}[(a),beginpenalty=10000]
\item Construct a one-hidden-layer ReLU network
\[
 F:\R^{|Q|+|\Sigma|}\to\R^{|Q|}
\]
with at most \(|Q||\Sigma|\) hidden units and an affine output layer such
that
\[
 F(e_q,e_a)=e_{\delta(q,a)}
 \qquad(q\in Q,\ a\in\Sigma).
\]
Give all hidden-unit and output formulas.
\item For an input string \(a_1\cdots a_T\), define
\[
 h_0=e_{q_0},
 \qquad
 h_t=F(h_{t-1},e_{a_t}),\qquad t\in[T].
\]
Prove that \(h_t\) is the one-hot encoding of the automaton state after
reading \(a_1\cdots a_t\).
\item Compare the recurrent implementation with a feedforward network obtained
by unrolling \(T\) copies of \(F\). State separately how the number of distinct
parameters, the number of executed cell evaluations, and the unrolled depth
depend on \(T\).
\end{enumerate}
\begin{hints}
On bit inputs, one ReLU can indicate that two specified coordinates are both
one. If a second hint is needed, introduce one hidden unit for every pair
\((q,a)\in Q\times\Sigma\), then collect the pairs that lead to the same next
state.
\end{hints}
\end{exercisechunk}
